\section{The coefficient and the positive phase}\label{sec:coefficients}

We first compute the coefficient for \eqref{eq:complex}, allowing an
arbitrary rectangular differential $D$. Adjoint symbols always refer
to the current Hermitian metrics. Put
\begin{equation}\label{eq:BCLE}
 B=D^\dagger D,\quad L=DD^\dagger,\quad
 C=\sum_j A_j^\dagger A_j,\quad E=DCD^\dagger,
 \quad P=e^{-B},\quad Q=e^{-L}.
\end{equation}
For a nonnegative Hermitian matrix $B$, define
\begin{align}
 R(B)&=\int_0^1e^{-sB}\,ds,\label{eq:R}\\
 K_B(C)&=\int_0^1e^{-(1-s)B}Ce^{-sB}\,ds,\label{eq:K}\\
 J_L(E)&=\int_{\substack{x,z\geq0\\x+z\leq1}}
                 e^{-xL}Ee^{-zL}\,dx\,dz.\label{eq:J}
\end{align}
In particular $R(0)=1$; no inverse of $B$ is implicit in this notation.

\begin{proposition}\label{prop:coefficient}
The superconnection coefficient of \cite[\S5.4]{hkkp26} on
\eqref{eq:complex}, including its degree sign, is
\begin{equation}\label{eq:coeff}
 \begin{split}
 \Omega_V&=z_0P+K_B(C),\\
 \Omega_U&=J_L(E)-z_0Q,\\
 \Omega_W&=z_1\id_W-\sum_j A_jR(B)A_j^\dagger.
 \end{split}
\end{equation}
If $n=\dim V$, $r=\dim U$ and $m_0=\dim W$, then
\begin{equation}\label{eq:central}
 \tr\Omega=(n-r)z_0+m_0z_1.
\end{equation}
Moreover,
\begin{equation}\label{eq:imaginary}
 \Ima\Omega_V=(\Ima z_0)P,\qquad
 \Ima(-\Omega_U)=(\Ima z_0)Q,\qquad
 \Ima\Omega_W=(\Ima z_1)\id_W .
\end{equation}
\end{proposition}

\begin{proof}
We specify the graded conventions to fix the signs. In the quiver
algebra of \cite[Definition~5.13]{hkkp26}, an arrow $\alpha:s\to t$
and its adjoint have degrees $1$ and $-1$, while the vertex elements
$e_i,f_i$ have degree zero. The nontrivial two-arrow products are
$\alpha\alpha^*=\ii f_t$ and
$\alpha^*\alpha=-\ii f_s$ for a matching arrow; cross-arrow products
vanish. Products involving an $f_i$ and an arrow, and products of two
$f_i$'s, vanish. Tensor multiplication has the Koszul sign
\[
 (x\otimes T)(y\otimes S)=(-1)^{|T||y|}xy\otimes TS.
\]
The involution on an even endomorphism is the ordinary adjoint, and
on an odd endomorphism it is $\ii$ times that adjoint.
The source superconnection and curvature are
\[
 \mathsf A=\sum_i e_i\otimes d_i+\sum_j\alpha_j\otimes A_j,\qquad
 F=(\mathsf A-\mathsf A^*)^2+\ii\sum_i f_i\otimes z_i.
\]
The coefficient is
$\Omega=(-1)^{\deg}\exp(-\ii F)_f$, where the subscript means
coefficient extraction after exponentiation.

Write $X=-\ii F$. Its vertex part is $-B$ on $V$, $-L$ on $U$
and zero on $W$. Its $f$ coefficients before exponentiation are
$C+z_0\id$ on $V$, $z_0\id$ on $U$, and
$z_1\id-\sum_j A_jA_j^\dagger$ on $W$.
Its mixed arrow terms are $-A_jD^\dagger:U\to W$ and
$-\ii DA_j^\dagger:W\to U$, with their respective algebra factors.
Their paired $f$ contribution is $-DCD^\dagger$ on $U$ and
$\sum_j A_jD^\dagger DA_j^\dagger$ on $W$.

The Duhamel expansion contributing to an $f$ coefficient terminates
at one $f$ insertion or two arrow insertions. There are no mixed
paths on $V$, giving $K_B(C)+z_0P$. On $U$, the two-insertion term
is $-J_L(E)$ before multiplication by the degree-one sign.
On $W$ the two-insertion contribution is
\[
 \sum_j\int_0^1(1-s)A_jD^\dagger e^{-sL}DA_j^\dagger\,ds.
\]
Using $D^\dagger e^{-sL}D=Be^{-sB}$ and
\[
 -\id+\int_0^1(1-s)Be^{-sB}\,ds=-R(B)
\]
proves \eqref{eq:coeff}.

The substitutions $s\mapsto1-s$ and $(x,z)\mapsto(z,x)$ show that
$K_B(C)$ and $J_L(E)$ are Hermitian. This proves
\eqref{eq:imaginary}; neither integral operator is asserted to preserve
positivity. Cyclicity of trace gives
\[
 \tr K_B(C)=\tr(CP),\qquad
 \tr J_L(E)
 =\tr\left(C\int_0^1sBe^{-sB}\,ds\right)
 =\tr(C(R(B)-P)).
\]
The arrow traces cancel. The nonzero spectra of $B$ and $L$ agree,
so $\tr P-\tr Q=n-r$, proving \eqref{eq:central}.
\end{proof}

\begin{lemma}\label{lem:phase}
If $\Ima T>0$ and $T=U_T|T|$ is its polar decomposition, then
every eigenvalue of $U_T$ lies in the open upper semicircle.
On a block of degree $q$, the matrix
$\Phi=\Arg U_T-\pi q\id$ is the unique Hermitian matrix with spectrum
in $(-\pi q,\pi(1-q))$ such that
$(-1)^qT e^{-\ii\Phi}$ is positive definite.
\end{lemma}

\begin{proof}
For a unit eigenvector $v$ of $U_T$ with eigenvalue $e^{\ii\theta}$,
\[
 0<\inner{v}{\Ima(T)v}
   =\sin\theta\,\inner{v}{|T|v}.
\]
Thus $0<\theta<\pi$. With the stated branch,
$(-1)^qT e^{-\ii\Phi}=TU_T^\dagger=(TT^\dagger)^{1/2}>0$.
Conversely positivity identifies $e^{\ii\Phi}$ as the polar unitary
of $(-1)^qT$. Uniqueness of polar decomposition and of the branch
give uniqueness of $\Phi$.
\end{proof}

For the trace functional $\mathcal L(a)=\tr(\Omega a)$,
the positivity requirement
$\mathcal L(e^{-\ii\Phi}a^\dagger a)\geq0$ for all $a$
is exactly the positivity of $\Omega e^{-\ii\Phi}$.
Lemma~\ref{lem:phase} therefore identifies the phase required by
\cite[\S4.2]{hkkp26} without a normality hypothesis.
The degree convention also respects shifts: the shift $[1]$ moves
degree $q$ to $q-1$ and increases the phase by $\pi$.
Orthogonal direct sums give direct sums of coefficients and phases,
and additive mass.

\begin{example}\label{ex:nonnormal}
Take $V=U=W=\CC^2$, $z_0=z_1=\ii$,
$D=\diag(1,\sqrt2)$ and
$A=\left(\begin{smallmatrix}1&1\\0&1\end{smallmatrix}\right)$.
Writing $p=e^{-1}$ and $q=e^{-2}$, formula \eqref{eq:coeff} gives
\[
 \Omega_V=
 \begin{pmatrix}(1+\ii)p&p-q\\p-q&(2+\ii)q\end{pmatrix},\qquad
 [\Omega_V,\Omega_V^\dagger]
 =2\ii(p-q)^2\begin{pmatrix}0&1\\-1&0\end{pmatrix}\ne0.
\]
Thus coefficient extraction does not preserve the normality of an
element in the graded algebra. This example contradicts the
normality clause of \cite[Lemma~5.16]{hkkp26}, while retaining the
strict imaginary-part property and the central-charge trace.
\end{example}

\subsection{Variation and global existence}
At an instant choose orthonormal frames. Along the metric phase flow,
the fixed maps have moving-frame derivatives
\begin{equation}\label{eq:mapflow}
 \dot D=\tfrac12(D\Phi_V-\Phi_UD),\qquad
 \dot A_j=\tfrac12(A_j\Phi_V-\Phi_WA_j).
\end{equation}
For example, a variation $\xi_q=h_q^{-1}\delta h_q$ gives
$\delta D=(\xi_UD-D\xi_V)/2$ and
$\delta A_j=(\xi_WA_j-A_j\xi_V)/2$ in such frames.
Substituting $\xi=-\Phi$ gives \eqref{eq:mapflow}.

All coefficient variations can be computed without denominators involving
eigenvalue differences. Define block matrices
\[
 N_V=\begin{pmatrix}-B&C\\0&-B\end{pmatrix},\quad
 N_U=\begin{pmatrix}-L&E&0\\0&0&\id\\0&0&-L\end{pmatrix},\quad
 N_R=\begin{pmatrix}-B&\id\\0&0\end{pmatrix}.
\]
Then $(e^{N_V})_{11}=P$, $(e^{N_V})_{12}=K_B(C)$,
$(e^{N_U})_{11}=Q$, $(e^{N_U})_{13}=J_L(E)$ and
$(e^{N_R})_{12}=R(B)$. The Fr\'echet derivative
\begin{equation}\label{eq:duhamel}
 \mathrm D\exp_N(\dot N)=\int_0^1e^{(1-s)N}\dot N e^{sN}\,ds
\end{equation}
and the product rule therefore determine every $\dot\Omega_q$.
In particular,
\[
 \dot\Omega_W=-\sum_j\bigl(\dot A_jRA_j^\dagger+
 A_j\dot R A_j^\dagger+A_jR\dot A_j^\dagger\bigr).
\]

\begin{lemma}\label{lem:massvariation}
Let $\Omega_q=U_q|\Omega_q|$ on each homogeneous block. Then
\begin{equation}\label{eq:massvariation}
 \dot m=\sum_q\Rea\tr(U_q^\dagger\dot\Omega_q).
\end{equation}
The coefficients and phase define a smooth metric vector field, whose
solutions exist for all $t\geq0$.
\end{lemma}

\begin{proof}
For an invertible matrix $\Omega=UH$, differentiation of
$H^2=\Omega^\dagger\Omega$ and taking trace after multiplication by
$H^{-1}$ gives
$\tr\dot H=\Rea\tr(U^\dagger\dot\Omega)$.
This proves \eqref{eq:massvariation}.
The polar decomposition and its upper-semicircle logarithm are smooth
on the domain in Lemma~\ref{lem:phase}.
If a phase derivative is needed, put
$S=U^\dagger\dot U$ and $Z=U^\dagger\dot\Omega$; it is determined by
\[
 HS+SH=Z-Z^\dagger,\qquad S^\dagger=-S,
\]
followed by differentiation of the chosen unitary logarithm.
The Sylvester operator is invertible because $H>0$.

The branches give $0<\Phi_V,\Phi_W<\pi\id$ and
$-\pi\id<\Phi_U<0$. In fixed coordinates the metric forms satisfy
\[
 e^{-\pi t}h_V(0)\leq h_V(t)\leq h_V(0),\qquad
 h_U(0)\leq h_U(t)\leq e^{\pi t}h_U(0),
\]
and the same degree-zero bounds hold on $W$.
On every finite time interval these inequalities confine the metrics
to a compact subset of the positive cone. The smooth ordinary
differential equation therefore extends past every finite time.
\end{proof}
